\section{System Architecture}
\label{sec:architecture}

% Facts verified 2026-06-10 against: mamai/config/{app_config.json,
% runtime_config.json, rag_assets.lock.json}, mamai/CLAUDE.md,
% config/prompts/system_en.txt, RagPipeline.kt.

MAM-AI is an Android application that answers clinical questions from
nurse-midwives entirely on-device: retrieval, generation, and all model
weights live on the phone, and no query ever leaves it. This design is
forced by the deployment context---health facilities in Zanzibar cannot
assume reliable internet connectivity---and it has a second benefit that
matters in a clinical setting: patient-related queries are never
transmitted to a third-party service. The cost is that every component
must fit the memory, compute, and thermal budget of a commodity
smartphone. This section describes the deployed configuration; the
corpus it retrieves from is described in Section~\ref{sec:corpus}, and
its empirical behaviour in Sections~\ref{sec:retrieval-results}
through~\ref{sec:latency}.

The application has three layers: a Flutter user interface, a Kotlin
native layer hosting the RAG pipeline, and two inference runtimes
(LiteRT-LM for generation, LiteRT/TFLite for embeddings) backed by a
SQLite vector store. All runtime hyperparameters are defined in a single
versioned configuration file (\texttt{runtime\_config.json}) that is
shared verbatim between the shipped application and the evaluation
harness, so that evaluation results refer to exactly the deployed
configuration (Section~\ref{sec:methodology}).

\subsection{Flutter UI and Platform Channels}

The user-facing layer is a Flutter application offering a single search
interface with suggestion chips and Markdown rendering of responses.
Flutter communicates with the native layer over two platform channels: a
\emph{method channel} through which the UI triggers initialization
(\texttt{ensureInit}) and submits queries (\texttt{generateResponse}),
and an \emph{event channel} over which the native layer streams results
back---first the retrieved documents, then generated tokens as they are
decoded. Streaming matters at on-device speeds: total generation takes
tens of seconds on target hardware (Section~\ref{sec:latency}), and
token-by-token display is the difference between a usable tool and an
apparently frozen one.

\subsection{The RAG Pipeline}

The native layer (\texttt{RagPipeline.kt}) manages three components: the
LLM runtime, the embedding model, and the vector store. Initialization
is lazy and asynchronous---models are loaded on first access, and
callers block on \texttt{awaitLlmReady()}, which propagates
initialization failures rather than silently unblocking. A companion
stream manager (\texttt{RagStream.kt}) enforces a single-query-at-a-time
concurrency constraint: duplicate submissions of the same prompt are
deduplicated, and a new query cancels the previous generation job before
starting. On a device where one generation saturates the SoC, admitting
concurrent queries would only slow both.

A query flows through the pipeline as follows: the query text is
embedded by Gecko; the vector store returns the top-3 chunks by cosine
similarity; the chunks, the system prompt, and the question are
assembled into a single prompt; and LiteRT-LM decodes a streaming
response.

\subsection{Gemma 4 E4B via LiteRT-LM}

Generation uses Gemma~4~E4B\footnote{TODO: cite the Gemma 4 technical
report.} in the int4-quantized \texttt{.litertlm} format (3.65\,GB),
executed by LiteRT-LM~0.11.0.\footnote{\url{https://github.com/google-ai-edge/LiteRT-LM}}
The engine is configured with a 4{,}096-token total context (prompt plus
response); this ceiling is not arbitrary but a deliberate safety margin
below a precision-induced failure mode at longer contexts that we
characterize in Section~\ref{sec:latency}. Sampling uses the model's
recommended deployment parameters (temperature 1.0, top-$p$ 0.95,
top-$k$ 64), with no explicit output-length cap.

The LLM backend defaults to CPU (XNNPACK) in production builds. GPU
execution (OpenCL; FP16 by default) is a compile-time opt-in, and if GPU
initialization fails at runtime the pipeline falls back to CPU
automatically. GPU decoding is 2--3.5$\times$ faster on supported
hardware (Section~\ref{sec:latency}), but the conservative default
reflects the fleet heterogeneity of the target market, where GPU driver
reliability cannot be assumed. A multi-token-prediction (speculative
decoding) mode was integrated and smoke-tested but is disabled by
default: on long retrieval-augmented prompts it measured 10--20\%
\emph{slower} than standard decoding.

\subsection{Gecko Embeddings and SQLite Vector Store}

Query embedding uses Gecko-110M~\cite{lee2024gecko} as a quantized
TFLite artifact (\texttt{Gecko\_1024\_quant.tflite}, 768-dimensional
output) with a SentencePiece tokenizer, running on CPU. Document
embeddings are \emph{not} computed on-device: the corpus is embedded
ahead of time by the producer pipeline (Section~\ref{sec:corpus}) and
shipped as a $\sim$260\,MB SQLite database of 63{,}650 pre-computed
vectors, stored in a compact binary format (a magic prefix followed by
768 little-endian \texttt{float32} values per row). The vector store is
the SQLite-backed semantic memory from Google AI Edge's local-agents RAG
library, queried with cosine similarity.

Each stored chunk carries a structured header
(\texttt{[SOURCE:$\langle$stem$\rangle$|PAGE:$\langle$n$\rangle$|CID:$\langle$hex$\rangle$]})
identifying its source document, page, and a stable 64-bit content ID.
The pipeline parses this header at retrieval time, which is what allows
the UI to display per-chunk citations and the evaluation harness to
trace any retrieved passage back to the exact guideline page that
produced it.

\subsection{Prompt Template and Retrieval Parameters}

Retrieval is fixed at top-3 with no similarity threshold (threshold
0.0): with a 4{,}096-token context shared between prompt and response,
three chunks of up to 2{,}500 characters are what comfortably fits
alongside the system prompt, the question, and room for the answer.

The system prompt (shipped in English and Swahili) defines the assistant
as clinical decision support for government nurse-midwives---explicitly
not specialist midwives---working in resource-limited facilities. Its
load-bearing instructions, each of which traces to an observed
failure mode or deployment requirement: simple short sentences for
second-language English speakers; responses under 200 words unless a
procedure requires more; bracketed citation markers ([1], [2], [3])
tied to the retrieved chunks; a fixed list of emergency presentations
(postpartum haemorrhage, eclampsia, cord prolapse, shoulder dystocia,
neonatal sepsis, among others) that must trigger immediate escalation
advice; a prohibition on recommending drug doses unless the retrieved
context explicitly states them; and an instruction to admit uncertainty
rather than guess. Retrieved chunks are injected under a labelled
context block, and the model is told to fall back to established
medical knowledge---and say so---when the retrieved context is
irrelevant. The behavioural consequences of this
safety-conservative prompt are measurable and discussed in
Section~\ref{sec:endtoend}.

\subsection{Model Distribution and Storage}

The APK ships without weights. On first launch the application downloads
roughly 4.5\,GB of artifacts: the LLM and embedder from public
HuggingFace repositories (\texttt{litert-community}), and the RAG bundle
from the producer repository's GitHub releases. The bundle is pinned by
a lock file (\texttt{rag\_assets.lock.json}) recording the bundle
version (v0.2.0 for this report), its download URL, the SHA-256 of its
manifest, and the producer commit that built it---so a given APK build
retrieves from a bit-exact, auditable corpus. All artifacts are cached
in application-external storage; after the initial download the
application runs fully offline. The first-launch download is the main
deployment friction in a low-connectivity setting, and
offline distribution (pre-loaded devices or local file transfer) is an
open operational item discussed in Section~\ref{sec:limitations}.
